\section{Aggregate geometry for the capped adjacency matrix}
\label{sec:sparse-extension}

Let $Z_{ik}=\ind{z_i=k}$ be the community indicator matrix.
The completed matrix $ZPZ^\top$ includes its diagonal, whereas
$\E A=ZPZ^\top-\diag(P_{z_i z_i})$. Their difference has norm at
most $p$ and is retained in the bounds below. The completed row
sums $\sum_jP_{z_i z_j}$ are block constant and comparable to $np$.
In this appendix put $\tau=b_0\bar D$ and
$G=\diag\{\max(D_i,\tau)\}$, with diagonal entries $G_i$.
Then $W=\tau^{1/2}G^{-1/2}$ is the algorithm's high-degree cap,
with $b_0$ sufficiently large for the structural class. Thus
$\widetilde A=WAW$ leaves every edge between vertices of degree at most $\tau$
unchanged. For analysis only, write
$Q=\widetilde A/\tau=G^{-1/2}AG^{-1/2}$ and
$\norm v_G^2=v^\top Gv$. The scalar rescaling $Q$ has the same
eigenvectors as the encoded adjacency matrix $\widetilde A$.

\begin{proposition}[Aggregate capped-adjacency geometry]\label{prop:sparse-floor}
Under \eqref{eq:structural}, let $np\to\infty$ and $np^2\le C_1$.
For every fixed $r>0$, with probability at least $1-C_rn^{-r}$,
$c{np}\le\tau\le C{np}$, $\tau\ge\max_i\sum_jP_{z_i z_j}$,
$c{n^2p}\le\sum_iD_i$, and
\begin{align}
 \norm{W(A-{ZPZ^\top})W}_{\op}&\le C_r\sqrt {np},
 &\sum_iG_i&\le C_r{n^2p},\label{eq:floor-noise}\\
 \sum_i\frac{(D_i-\sum_jP_{z_i z_j})^2}{G_i}
 +\sum_i\frac{(G_i-\tau)^2}{G_i}
 &\le C_rn.\label{eq:floor-degree-energy}
\end{align}
In particular,
$\norm{G^{-1/2}(A-{ZPZ^\top})G^{-1/2}}_{\op}\le C_r/\sqrt {np}$.
Let $U$ contain the $K$ eigenvectors of $\widetilde A$ with largest absolute
eigenvalues, and put $\Pi=UU^\top$ and $H=G^{-1/2}\Pi G^{1/2}$.
The eigenvalue equation gives $\widehat A=AH$ for the reconstruction
$\widehat A=W^{-1}U\Lambda U^\top W^{-1}$ in \cref{restore:algorithm}.
The retained eigenvalues of $\widetilde A$ have magnitude at least $c_r{np}$;
the remaining eigenvalues have magnitude at most $C_r\sqrt {np}$.
There is an orthogonal matrix $O$ such that the orthonormal population
basis
\[
 V=G^{-1/2}Z(Z^\top G^{-1}Z)^{-1/2}O
\]
satisfies
\[
 \norm{U-V}_F\le C_r/\sqrt {np},
 \qquad \sqrt{G_i}\norm{V_{i\cdot}}_2\le C_r\sqrt{p}.
\]
Moreover,
\begin{equation}\label{eq:floor-block-residual}
 \norm{(I-\Pi)G^{1/2}Z}_F\le C_r\sqrt n,
\end{equation}
and every block-constant $f$ with $\norm f_\infty\le C_f$ obeys
\begin{equation}\label{eq:floor-weighted-errors}
 \norm{(I-H)f}_G\le C_rC_f\sqrt n,\qquad
 \norm{(I-H)f}_2^2\le \frac{C_rC_f^2}{p}.
\end{equation}
The same assertions hold simultaneously after deletion of any fixed
number of vertices, with each cap recomputed from the remaining graph,
after increasing the probability parameter. No raw minimum-degree
condition is imposed.
\end{proposition}

\begin{proof}
The Bernoulli transform of the total number of edges gives
$c{n^2p}\le\sum_iD_i\le C{n^2p}$ except with probability $e^{-c{n^2p}}$.
For large enough $b_0$, this event implies
\[
 c{np}\le\tau=b_0\bar D\le C{np},
 \qquad \tau\ge\max_i\sum_jP_{z_i z_j}.
\]
Since $\E A_{ij}\le p$, put
$W_{h,ii}=\sqrt{\min(1,2{np}/D_i)}$, with value one at zero degrees.
The changed set $\mathcal S=\{i:D_i>2{np}\}$ contains at most
$\frac{10}{p}$ vertices except with probability $e^{-cn}$. To check the
strength of this bound, set $m=\lfloor\frac{10}{p}\rfloor+1$. If $m\le n$
and $m$ vertices exceed the threshold, some set $T$ of size $m$ has
$\sum_{i\in T}D_i>2m{np}$. This is an independent-edge sum with
coefficients in $\{1,2\}$ and mean at most $m{np}$. At
$t=(\log2)/2$, its exponential transform bounds the tail by
\[
 \exp\{-(\log2-1/2)m{np}\},
\]
because $(e^{tc}-1)/c\le(e^{2t}-1)/2=1/2$ for $c=1,2$.
Since $\log\binom nm\le m\log(en/m)=o(m{np})$, a union over $T$
gives the stated bound. If $m>n$, the claim is automatic.

Every row of $W_hAW_h$ has squared Euclidean norm at most
$W_{h,ii}^2D_i\le2{np}$. The row-norm form of the regularized adjacency
theorem of \citet[Theorem 2.1 and Remark 2.2]{llv2017} therefore yields
\[
 \norm{W_hAW_h-\E A}_{\op}\le C_r\sqrt {np}
\]
with failure probability at most $C_rn^{-r}$. The theorem permits
these graph-dependent weights. Also
\[
 \norm{W_h{ZPZ^\top} W_h-{ZPZ^\top}}_{\op}
 \le Cp\sqrt{n|\mathcal S|}\le C\sqrt {np},
\]
by the Frobenius norm: the difference is supported on rows or columns
in $\mathcal S$. The completed-mean diagonal has norm $O(p)$.
Consequently $\norm{W_h(A-{ZPZ^\top})W_h}_{\op}\le C_r\sqrt {np}$.
The exact diagonal comparison is
\[
 G^{-1/2}=(2{np})^{-1/2}
 \diag\!\left\{\sqrt{\frac{\max(D_i,2{np})}{G_i}}\right\}W_h.
\]
The displayed diagonal multiplier is bounded, since
$G_i\ge\max(D_i,c{np})$.
Conjugation gives the stated bound for the noise in $Q$;
multiplication by $\tau\asymp {np}$ proves the adjacency-scale bound
in \eqref{eq:floor-noise}. Since $G_i\le D_i+\tau$, the total
weight $\sum_iG_i$ is at most $C{n^2p}$.

Applying the scalar-rescaled noise to $G^{1/2}\one$ gives
\[
 \norm{G^{-1/2}(A-{ZPZ^\top})\one}_2
 \le C_r{(np)}^{-1/2}\norm{G^{1/2}\one}_2\le C_r\sqrt n,
\]
which is the first sum in \eqref{eq:floor-degree-energy}.
For the second sum, $G_i-\tau=0$ when $D_i\le\tau$; otherwise
$0\le G_i-\tau=D_i-\tau\le D_i-\sum_jP_{z_i z_j}$. Thus, coordinatewise,
\[
 \frac{(G_i-\tau)^2}{G_i}
 \le\frac{(D_i-\sum_jP_{z_i z_j})^2}{G_i}.
\]
The same bound also gives
$\sum_i(1-W_{ii})^2\le \frac{C_r}{p}$, since
$(1-W_{ii})^2\le(G_i-\tau)^2/(\tau G_i)$.

Let $J=Z^\top G^{-1}Z$. For every block $k$, Cauchy--Schwarz,
comparable group sizes, $G_i\ge c{np}$ and the total-degree bound give
\[
 \frac{c}{p}\le
 \frac{n_k^2}{\sum_{i:z_i=k}G_i}
 \le J_{kk}=\sum_{i:z_i=k}G_i^{-1}\le\frac{C}{p}.
\]
The nonzero spectrum of the capped population matrix $W{ZPZ^\top} W$ is
the spectrum of $\tau J^{1/2}PJ^{1/2}$. The bound
$\sigma_K(P)\ge c p$ gives a separation of order ${np}$ from zero.
The adjacency noise bound and
subspace perturbation then give the stated eigenspace and eigenvalue
bounds. The formula for $V$
and the lower bound on $J$ give its weighted row bound.

The matrix $\tau G^{-1/2}Z$ belongs to the population signal space
$\col(WZ)=\col(G^{-1/2}Z)$. Furthermore,
\[
 \norm{G^{1/2}Z-\tau G^{-1/2}Z}_F^2
 =\sum_i\frac{(G_i-\tau)^2}{G_i}\le C_rn,
 \qquad \norm{G^{1/2}Z}_F^2\le C_r{n^2p}.
\]
Together with the projector error $O({(np)}^{-1/2})$, these imply
\eqref{eq:floor-block-residual}. The matrix $H$ is the
$G$-orthogonal projector corresponding to $\Pi$, so the first bound
in \eqref{eq:floor-weighted-errors} follows by applying the residual bound to the $K$ block values of $f$.
The second uses $G_i\ge c{np}$.

Deleting a fixed number of vertices preserves the structural bounds.
Apply the same argument to each recomputed graph and choose a larger
fixed probability parameter before the finite-depth union bound.
\end{proof}
